\chapter{操作数相关}
Python是弱类型语言, 变量的类型是根据赋值的值来决定的。

\section{概述}
变量是数据操作和数值计算的起始点，
算法用到的各种数据都是存储在变量内。

Python 是一门"弱类型"语言，所有的变量无须声明即可使用，变量的数据类型可以随时改变. 
Python3 中有六个标准的数据类型，分类如下:
\begin{enumerate}[label=\arabic*)]
    \item 不可变数据类型
    \begin{enumerate}[itemindent=1em,label=\arabic*)]
        \item Number (数字或数值)
        \item String (字符串)
        \item Thple(元组)
    \end{enumerate}
    
    
    \item 可变数据类型——序列
    \begin{enumerate}[itemindent=1em,label=\arabic*)]
        \item List (列表)
        \item Set (集合)
        \item Dictionary (字典)
    \end{enumerate}
\end{enumerate}

\subsection{标识符(变量名)}
\begin{enumerate}[label=\arabic*)]
    \item 第一个字符必须以字母(\codeinline{a-z}, \codeinline{A-Z})或下划线 \codeinline{\_} 。
    \item 标识符的其他的部分由字母、数字和下划线组成。
    \item 标识符对大小写敏感，count 和 Count 是不同的标识符。
    \item 标识符对长度无硬性限制，但建议保持简洁（一般不超过 20 个字符）。
    \item 禁止使用保留关键字，如 \codeinline{if}、\codeinline{for}、\codeinline{class} 等不能作为标识符。
\end{enumerate}
\begin{lstlisting}[style=demostyle,numbers=none]
# 合法标识符：    
age = 25
user_name = "Alice"
_total = 100
MAX_SIZE = 1024
calculate_area()
StudentInfo
__private_var

# 非法标识符：
2nd_place = "silver"    # 错误：以数字开头
user-name = "Bob"       # 错误：包含连字符
class = "Math"          # 错误：使用关键字
$price = 9.99          # 错误：包含特殊字符
for = "loop"           # 错误：使用关键字
\end{lstlisting}

Python 3 允许使用 Unicode 字符作为标识符，可以用中文作为变量名，非 ASCII 标识符也是允许的了。
\begin{lstlisting}[style=demostyle,numbers=none]
姓名 = "张三"  # 合法
π = 3.14159   # 合法
\end{lstlisting}























